\chapter{Renormalization of the Quark Meson model}

\section{Renormalization of the Quark--Meson Model}

\subsection{Scope and approximation}

In this subsection we consider the renormalization of the two-flavor
Quark--Meson model in the vacuum limit, corresponding to vanishing
temperature and chemical potentials,
\begin{equation}
T = 0, \qquad \mu_u = \mu_d = 0 .
\end{equation}
The aim is to determine how quantum fluctuations modify the parameters
of the classical Lagrangian and how the ultraviolet divergences
generated by loop diagrams can be absorbed into redefinitions of the
parameters of the model.

The analysis is performed at one-loop order in perturbation theory.
At this order the quantum corrections to Green's functions arise from
diagrams containing a single closed loop of internal propagators
\cite{PeskinSchroeder1995,WeinbergQFT}. In the Quark--Meson model both
quark and meson fields can in principle propagate in such loops.
However, the relative importance of different diagrams can be
organized using the large-$N_c$ expansion, where $N_c$ denotes the
number of colors.

In the large-$N_c$ limit of QCD, diagrams containing closed quark
loops scale proportionally to $N_c$, whereas diagrams involving only
mesonic propagators scale as $\mathcal{O}(1)$
\cite{SchaeferWambach2008}. Quark-loop contributions therefore
dominate the one-loop corrections, while mesonic loop corrections are
suppressed by powers of $1/N_c$.

Throughout this section we adopt this approximation and retain only
the loop diagrams generated by quark fields. Mesonic loop corrections
to propagators and vertices are neglected. This is the same
approximation that was used in the previous section when deriving the
mean-field thermodynamic potential, where the quark fields were
integrated out while the mesonic fields were treated as classical
background fields \cite{Folkestad2018,Andersen2024}. Renormalizing the
vacuum theory within the same approximation ensures that the
parameters of the model remain consistent with the thermodynamic
analysis already performed.

\subsection{Renormalized Lagrangian and counterterms}

Renormalization is implemented by distinguishing between the
\emph{bare} quantities that appear in the original Lagrangian and the
\emph{renormalized} quantities that are identified with physical
parameters of the theory. The bare quantities contain ultraviolet
divergences generated by loop diagrams, while the renormalized
parameters remain finite after these divergences have been absorbed
into suitable counterterms \cite{PeskinSchroeder1995,WeinbergQFT}.

We therefore rewrite the Quark--Meson Lagrangian in terms of bare
fields and bare parameters. For the mesonic sector this takes the
form
\begin{equation}
\mathcal{L}_{\mathrm{mes}}^{(0)}
=
\frac{1}{2}(\partial_\mu \sigma_0)^2
+
\frac{1}{2}(\partial_\mu \boldsymbol{\pi}_0)^2
-
\frac{m_0^2}{2}(\sigma_0^2+\boldsymbol{\pi}_0^2)
-
\frac{\lambda_0}{4}(\sigma_0^2+\boldsymbol{\pi}_0^2)^2
+
h_0 \sigma_0 ,
\end{equation}
where the subscript $0$ indicates bare quantities.

The bare fields are related to renormalized fields through
wavefunction renormalization factors,
\begin{equation}
\sigma_0 = Z_\sigma^{1/2}\sigma,
\qquad
\boldsymbol{\pi}_0 = Z_\pi^{1/2}\boldsymbol{\pi}.
\end{equation}
Similarly, the parameters of the mesonic potential are written as
\begin{align}
m_0^2 &= m^2 + \delta m^2, \\
\lambda_0 &= \lambda + \delta \lambda, \\
h_0 &= h + \delta h ,
\end{align}
where $m^2$, $\lambda$, and $h$ denote renormalized parameters and the
quantities $\delta m^2$, $\delta\lambda$, and $\delta h$ are
counterterms.

Substituting these relations into the bare Lagrangian allows the
theory to be written as the sum of a renormalized Lagrangian and a
counterterm contribution,
\begin{equation}
\mathcal{L}_{\mathrm{mes}}^{(0)}
=
\mathcal{L}_{\mathrm{mes}}
+
\mathcal{L}_{\mathrm{ct}} .
\end{equation}
The renormalized part $\mathcal{L}_{\mathrm{mes}}$ has the same form
as the classical Lagrangian written in terms of the renormalized
fields and parameters. The counterterm Lagrangian
$\mathcal{L}_{\mathrm{ct}}$ collects the additional terms generated by
the renormalization constants. These terms give rise to
counterterm vertices in perturbation theory.

In the present approximation only quark-loop diagrams contribute to
the quantum corrections. The ultraviolet divergences produced by
these loops are therefore absorbed into the counterterms introduced
above. Their explicit values will be determined later by imposing
renormalization conditions on the one- and two-point functions of the
mesonic fields.

\subsection{Expansion around the vacuum}

The vacuum structure of the model was analyzed earlier in the discussion
of the Quark-Meson Lagrangian, where the mesonic fields were decomposed
into spacetime-independent background fields and fluctuations. The
ground state was found to correspond to a nonvanishing expectation
value of the scalar field, while the pion expectation values vanish,
\begin{equation}
\langle \sigma \rangle = v,
\qquad
\langle \boldsymbol{\pi} \rangle = 0 .
\end{equation}

For the loop calculations considered here it is convenient to expand
the scalar field around this vacuum configuration. We therefore write
\begin{equation}
\sigma(x) = v + \tilde{\sigma}(x),
\end{equation}
where $v$ denotes the vacuum expectation value and
$\tilde{\sigma}(x)$ represents fluctuations about the vacuum.

Substituting this decomposition into the Yukawa interaction term of
the Lagrangian generates a fermion mass term proportional to
$\bar q q$. As discussed previously, this leads to the constituent
quark mass
\begin{equation}
M_q = g v .
\end{equation}
In perturbative calculations the quark propagator therefore involves
the mass $M_q$, while the fluctuation field $\tilde{\sigma}$ describes
propagating scalar excitations.

The expansion around the vacuum also determines the interaction
vertices relevant for loop diagrams. In particular, the Yukawa
interaction generates couplings between quarks and the mesonic
fluctuation fields $\tilde{\sigma}$ and $\boldsymbol{\pi}$. In
addition, the expansion of the mesonic potential produces terms
linear in $\tilde{\sigma}$. These terms correspond to tadpole
interactions and give rise to one-point diagrams in perturbation
theory. In the renormalized theory such contributions must cancel so
that the vacuum expectation value remains fixed at $v$.

\subsection{Diagrammatic structure at one loop}

Having established the expansion around the vacuum, we now identify
the loop diagrams that contribute to the renormalization of the
theory. As discussed in the previous subsection, the relevant quantum
corrections arise from quark loops, while mesonic loop diagrams are
neglected in the large-$N_c$ approximation. The renormalization
conditions will therefore be imposed on the one- and two-point
functions of the mesonic fields generated by quark loops
\cite{PeskinSchroeder1995,Folkestad2018}.

The simplest quantity to consider is the one-point function of the
scalar field. After shifting the field according to
$\sigma(x)=v+\tilde{\sigma}(x)$, the expansion of the mesonic
potential generates terms linear in the fluctuation field
$\tilde{\sigma}$. These terms correspond to tadpole interactions,
which produce one-point diagrams containing a single closed quark
loop. Such diagrams represent quantum corrections to the vacuum
expectation value of the scalar field. In the renormalized theory the
tadpole contributions must cancel so that the vacuum remains located
at $\langle\sigma\rangle=v$.

At the same order in perturbation theory, quark loops also generate
corrections to the two-point functions of the mesonic fields. These
corrections correspond to the self-energies of the pion and sigma
fields. The relevant diagrams consist of a quark loop with two
external meson lines attached through Yukawa vertices. For the pion
field this produces the pion self-energy $\Sigma_\pi(p^2)$, while the
corresponding diagram with scalar vertices generates the sigma
self-energy $\Sigma_\sigma(p^2)$.

In addition to the loop diagrams, perturbation theory also contains
diagrams generated by the counterterm Lagrangian. 
These counterterm diagrams have the same
topology as the loop diagrams but contain counterterm vertices rather
than internal propagators. Their role is to cancel the ultraviolet
divergences produced by the loop integrals so that the renormalized
Green's functions remain finite.

The renormalization program therefore involves three classes of
contributions: the quark-loop tadpole diagram, the quark-loop
corrections to the pion and sigma two-point functions, and the
corresponding counterterm diagrams. In the following subsections these
diagrams will be evaluated explicitly and the counterterms will be
fixed by imposing appropriate renormalization conditions.

\subsection{Tadpole renormalization}

We first consider the one-point function of the shifted scalar field
\(\tilde{\sigma}\). In the vacuum, the renormalized theory must still
be expanded around the physical minimum. This means that the
renormalized one-point function must vanish at the vacuum
configuration,
\begin{equation}
\Gamma^{(1)}_{\sigma}(v)=0 .
\label{eq:qm_tadpole_condition}
\end{equation}
Equivalently, the effective potential must satisfy the stationarity
condition
\begin{equation}
\left.
\frac{d\Omega_{\mathrm{eff}}}{d\sigma}
\right|_{\sigma=v}
=0 .
\label{eq:qm_stationarity_condition}
\end{equation}
This requirement ensures that the vacuum expectation value remains
equal to \(v\) after the one-loop corrections are included
\cite{Folkestad2018,Andersen2024}.

The corresponding tadpole cancellation structure is shown in
Fig.~\ref{fig:qm_tadpole_cancellation}.

\begin{figure}[H]
\centering
\begin{minipage}{0.4\textwidth}
\raggedleft
\feynmandiagram [layered layout, horizontal=b to c] {
    a [particle=\(\tilde{\sigma}\)] -- b 
    -- [fermion, half left] c
    -- [fermion, half left] b,
};
\end{minipage}
\begin{minipage}{0.1\textwidth}
\centering
\(
=
\)
\end{minipage}
\begin{minipage}{0.4\textwidth}
\raggedright
\feynmandiagram [layered layout, horizontal=a to b] {
    a [particle=\(\tilde{\sigma}\)] -- b [crossed dot]
};
\end{minipage}
\caption{Tadpole cancellation in the renormalized theory. The counterterm is chosen so that the renormalized one-point function vanishes and the vacuum remains fixed at \(v\).}
\label{fig:qm_tadpole_cancellation}
\end{figure}

At leading order in the large-\(N_c\) approximation, the only loop
contribution to the one-point function comes from the quark tadpole
diagram. After shifting the scalar field,
\(\sigma(x)=v+\tilde{\sigma}(x)\), the relevant Yukawa vertex is
\(-g\,\bar q \tilde{\sigma} q\), and the corresponding one-loop
contribution is proportional to the integral
\begin{equation}
A(M_q^2)
\equiv
\Lambda^{4-d}
\int \frac{d^d k}{(2\pi)^d}\,
\frac{1}{k^2-M_q^2+i0},
\qquad d=4-2\epsilon ,
\label{eq:qm_A_integral}
\end{equation}
where \(M_q \equiv gv\) is the constituent quark mass and \(\Lambda\)
is the renormalization scale introduced in dimensional
regularization \cite{Folkestad2018,Andersen2024}.

Including also the tadpole counterterm \(\delta t\), the one-point
function can be written as
\begin{equation}
i\Gamma^{(1)}_{\sigma}
=
-i\bigl(m_\pi^2 v - h\bigr)
- 8 N_c g M_q\, A(M_q^2)
+ i\,\delta t .
\label{eq:qm_one_point_function}
\end{equation}
The first term is the tree-level contribution. Since the classical
vacuum satisfies
\begin{equation}
m_\pi^2 v - h = 0 ,
\label{eq:qm_tree_level_stationarity}
\end{equation}
the one-point function at this order reduces to the sum of the quark
tadpole and the counterterm,
\begin{equation}
i\Gamma^{(1)}_{\sigma}
=
- 8 N_c g M_q\, A(M_q^2)
+ i\,\delta t .
\label{eq:qm_one_point_reduced}
\end{equation}

Imposing the renormalization condition
\eqref{eq:qm_tadpole_condition} therefore fixes the tadpole
counterterm to be
\begin{equation}
\delta t
=
-\,8 i N_c g M_q\, A(M_q^2)
=
-\,8 i N_c g^2 v\, A(M_q^2).
\label{eq:qm_tadpole_counterterm}
\end{equation}
The counterterm is thus chosen to cancel the full one-loop tadpole
contribution at the vacuum.

This result has a simple interpretation. Quantum fluctuations of the
quark fields tend to shift the position of the minimum of the
effective potential. The tadpole counterterm is introduced precisely
to remove this shift, so that the renormalized vacuum remains located
at the physical condensate \(v\). With the one-point function fixed in
this way, the remaining counterterms can be determined from the
renormalization of the pion and sigma two-point functions.

\subsection{Meson self-energies}

We now evaluate the quark-loop contributions to the mesonic two-point
functions. Since we work in the vacuum and retain only quark loops,
the relevant internal propagator is the fermion propagator with the
constituent quark mass,
\begin{equation}
S(k)=\frac{i(\slashed{k}+M_q)}{k^2-M_q^2+i0},
\qquad
M_q = gv .
\end{equation}
The pion and sigma self-energies are obtained by attaching two
external meson lines to a closed quark loop through the Yukawa
vertices. The calculation follows standard one-loop methods in
dimensional regularization \cite{PeskinSchroeder1995,WeinbergQFT} and
the explicit derivation given in \cite{Folkestad2018}. The same
one-particle-irreducible self-energies are also used in the on-shell
parameter fixing discussed in \cite{Andersen2024}.

The one-loop self-energy diagrams are displayed in
Fig.~\ref{fig:qm_selfenergy_diagrams}.

\begin{figure}[H]
\centering
\begin{minipage}{0.47\textwidth}
\centering
\feynmandiagram [layered layout, horizontal=b to c] {
    a [particle=\(\pi\)] -- b
      -- [fermion, half left] c
      -- [fermion, half left] b,
    c -- d [particle=\(\pi\)],
};
\end{minipage}
\hfill
\begin{minipage}{0.47\textwidth}
\centering
\feynmandiagram [layered layout, horizontal=b to c] {
    a [particle=\(\tilde{\sigma}\)] -- b
      -- [fermion, half left] c
      -- [fermion, half left] b,
    c -- d [particle=\(\tilde{\sigma}\)],
};
\end{minipage}
\caption{Quark-loop contributions to the pion and sigma self-energies at one loop. The pion diagram contains pseudoscalar Yukawa vertices, while the sigma diagram contains scalar Yukawa vertices.}
\label{fig:qm_selfenergy_diagrams}
\end{figure}

In the following we make use of the integral \(A(M_q^2)\) defined in 
Eq. \eqref{eq:qm_A_integral} and the standard scalar two-propagator integral
\begin{equation}
B(p^2)
\equiv
\Lambda^{4-d}
\int \frac{d^d k}{(2\pi)^d}\,
\frac{1}{\bigl[(k+p)^2-M_q^2+i0\bigr]\bigl[k^2-M_q^2+i0\bigr]},
\label{eq:qm_B_def}
\end{equation}
where \(d=4-2\epsilon\) and \(\Lambda\) is the renormalization scale
\cite{Folkestad2018,Andersen2024}.

We first consider the pion self-energy. The Yukawa interaction with
the pion field is
\begin{equation}
-ig\,\bar q \gamma^5 \tau_i \pi_i q ,
\end{equation}
so the one-loop pion self-energy is
\begin{equation}
i\Sigma_\pi(p^2)
=
(-1)\,2N_c\,g^2 \Lambda^{4-d}
\int\frac{d^d k}{(2\pi)^d}
\mathrm{Tr}\!\left[
\gamma^5
\frac{i(\slashed{k}+\slashed{p}+M_q)}{(k+p)^2-M_q^2+i0}
\gamma^5
\frac{i(\slashed{k}+M_q)}{k^2-M_q^2+i0}
\right].
\label{eq:qm_sigma_pi_start}
\end{equation}
The factor \((-1)\) comes from the closed fermion loop, while the
factor \(2N_c\) accounts for the two light flavors and the \(N_c\)
colors \cite{Folkestad2018}. Using
\(\{\gamma^5,\gamma^\mu\}=0\), \((\gamma^5)^2=1\), and
\(\mathrm{Tr}(\gamma^\mu\gamma^\nu)=4g^{\mu\nu}\), the Dirac trace becomes
\begin{equation}
\mathrm{Tr}\!\left[
\gamma^5(\slashed{k}+\slashed{p}+M_q)\gamma^5(\slashed{k}+M_q)
\right]
=
4\bigl(k^2+p\!\cdot\! k-M_q^2\bigr).
\end{equation}
After reducing the numerator against the denominators and expressing
the result in terms of \(A(M_q^2)\) and \(B(p^2)\), one finds
\begin{equation}
i\Sigma_\pi(p^2)
=
-8N_c g^2
\left[
A(M_q^2)-\frac{1}{2}p^2 B(p^2)
\right].
\label{eq:qm_pion_selfenergy_final}
\end{equation}
This result is the same for all three pion fields, as required by the
unbroken \(SU(2)_V\) symmetry of the vacuum \cite{Folkestad2018}.

We next consider the sigma self-energy. The corresponding Yukawa
interaction is
\begin{equation}
-g\,\bar q \tilde{\sigma} q ,
\end{equation}
which gives the one-particle-irreducible quark-loop contribution
\begin{equation}
i\Sigma_{\sigma}^{\mathrm{1PI}}(p^2)
=
(-1)\,2N_c\,(-ig)^2 \Lambda^{4-d}
\int\frac{d^d k}{(2\pi)^d}
\mathrm{Tr}\!\left[
\frac{i(\slashed{k}+\slashed{p}+M_q)}{(k+p)^2-M_q^2+i0}
\frac{i(\slashed{k}+M_q)}{k^2-M_q^2+i0}
\right].
\label{eq:qm_sigma_selfenergy_start}
\end{equation}
The Dirac trace is now
\begin{equation}
\mathrm{Tr}\!\left[
(\slashed{k}+\slashed{p}+M_q)(\slashed{k}+M_q)
\right]
=
4\bigl(k^2+p\!\cdot\! k+M_q^2\bigr),
\end{equation}
which differs from the pion case by the sign of the mass term. The
loop integral therefore evaluates to
\begin{equation}
i\Sigma_{\sigma}^{\mathrm{1PI}}(p^2)
=
-8N_c g^2
\left[
A(M_q^2)-\frac{1}{2}\bigl(p^2-4M_q^2\bigr)B(p^2)
\right].
\label{eq:qm_sigma_1PI_final}
\end{equation}
The difference between \eqref{eq:qm_pion_selfenergy_final} and
\eqref{eq:qm_sigma_1PI_final} reflects the different Dirac structure
of the scalar and pseudoscalar Yukawa vertices.

In the shifted theory the sigma two-point function also receives a
tadpole-induced contribution. It arises from the cubic
\(\tilde{\sigma}^3\) interaction in the mesonic sector combined with
the quark tadpole discussed in the previous subsection. At one-loop
order this contribution is represented by the tadpole-dressed sigma
two-point diagram shown in Fig.~\ref{fig:qm_sigma_tadpole_selfenergy},
and evaluates to

\begin{figure}[t]
\centering
\begin{tikzpicture}
    \begin{feynman}
        \vertex (a) {\(\tilde{\sigma}\)};
        \vertex[right=1.9cm of a] (b);
        \vertex[above=1.7cm of b] (d);
        \vertex[right=1.7cm of b] (c) {\(\tilde{\sigma}\)}; 
        \vertex[above=1.7cm of d] (e);
        \diagram* {
            (a) -- (b) -- (c),
            (b) -- (d),
            (d) -- [fermion, half left] (e),
            (e) -- [fermion, half left] (d),
        };
    \end{feynman}
\end{tikzpicture}
\caption{Tadpole-induced contribution to the sigma self-energy in the shifted theory. The lower line is the sigma two-point function, while the upper quark loop is attached through the cubic \(\tilde{\sigma}^3\) interaction generated after the field shift.}
\label{fig:qm_sigma_tadpole_selfenergy}
\end{figure}

\begin{equation}
i\Sigma_{\sigma}^{\mathrm{tad}}(p^2)
=
\frac{8N_c\lambda M_q^2}{m_\sigma^2}\,A(M_q^2),
\label{eq:qm_sigma_tadpole_part}
\end{equation}
so that the full sigma self-energy before renormalization is the sum
\begin{equation}
i\Sigma_\sigma(p^2)
=
i\Sigma_{\sigma}^{\mathrm{1PI}}(p^2)
+
i\Sigma_{\sigma}^{\mathrm{tad}}(p^2).
\label{eq:qm_sigma_full_before_ct}
\end{equation}
In the pion channel an analogous tadpole-induced contribution arises
from the $\tilde{\sigma}\pi^2$ interaction generated by the shift of
the scalar field. However, this term is proportional to the same
tadpole amplitude that is removed by the renormalization condition
$\Gamma_\sigma^{(1)}(v)=0$. Once the vacuum condition is imposed,
the corresponding pion contribution therefore vanishes together with
the sigma tadpole \cite{Folkestad2018}.

Equations \eqref{eq:qm_pion_selfenergy_final} and
\eqref{eq:qm_sigma_full_before_ct} give the unrenormalized one-loop
meson self-energies generated by quark fluctuations. They contain
ultraviolet divergences through the loop integrals \(A(M_q^2)\) and
\(B(p^2)\). In the next subsection these divergences are absorbed into
mass and wavefunction counterterms by imposing on-shell
renormalization conditions.

\subsection{On-shell renormalization conditions}

We now determine the counterterms associated with the mesonic
two-point functions. In the on-shell scheme the renormalized masses
are defined to coincide with the physical pole masses of the
propagators. The renormalized inverse propagator must therefore have
a zero at the physical mass and a residue equal to unity
\cite{PeskinSchroeder1995,WeinbergQFT,Folkestad2018,Andersen2024}.

After imposing the tadpole condition from the previous subsection,
the non-1PI pieces in the two-point functions are proportional to the
tadpole condition and therefore drop out. The inverse propagators may
then be written solely in terms of the 1PI self-energies as
\begin{equation}
G_\sigma^{-1}(p^2)
=
p^2 - m_\sigma^2
+ \Sigma_\sigma^{\mathrm{1PI}}(p^2)
+ \Sigma_{\sigma,\mathrm{ct}}^{\mathrm{1PI}}(p^2),
\label{eq:qm_inv_prop_sigma}
\end{equation}
\begin{equation}
G_\pi^{-1}(p^2)
=
p^2 - m_\pi^2
+ \Sigma_\pi^{\mathrm{1PI}}(p^2)
+ \Sigma_{\pi,\mathrm{ct}}^{\mathrm{1PI}}(p^2),
\label{eq:qm_inv_prop_pi}
\end{equation}
where \(m_\sigma\) and \(m_\pi\) denote the renormalized pole masses.
Here only the one-particle-irreducible self-energies appear
explicitly, since the tadpole condition fixes the vacuum at \(v\) and
eliminates the associated non-1PI contributions
\cite{Folkestad2018}.

The counterterm contribution to the two-point function has the
standard form
\begin{equation}
\Sigma_{a,\mathrm{ct}}^{\mathrm{1PI}}(p^2)
=
-\delta m_a^2
+
(p^2-m_a^2)\,\delta Z_a,
\qquad
a\in\{\pi,\sigma\},
\label{eq:qm_selfenergy_ct_general}
\end{equation}
where \(\delta m_a^2\) is the mass counterterm and \(\delta Z_a\) is
the wavefunction renormalization constant for the corresponding
mesonic field.

The first on-shell condition states that the full propagator must have
a pole at the physical mass. Equivalently, the inverse propagator must
vanish at \(p^2=m_a^2\),
\begin{equation}
G_a^{-1}(m_a^2)=0,
\qquad
a\in\{\pi,\sigma\}.
\label{eq:qm_OS_pole_condition}
\end{equation}
Using Eqs.~\eqref{eq:qm_inv_prop_sigma}--\eqref{eq:qm_selfenergy_ct_general},
this gives
\begin{equation}
\Sigma_a^{\mathrm{1PI}}(m_a^2)
+
\Sigma_{a,\mathrm{ct}}^{\mathrm{1PI}}(m_a^2)
=0.
\end{equation}
Since the term proportional to \(\delta Z_a\) vanishes at
\(p^2=m_a^2\), one obtains directly
\begin{equation}
\delta m_a^2
=
\Sigma_a^{\mathrm{1PI}}(m_a^2).
\label{eq:qm_delta_ma_general}
\end{equation}

Using the explicit one-loop self-energies from the previous
subsection, the mass counterterms become
\begin{equation}
\delta m_\sigma^2
=
8 i N_c g^2
\left[
A(M_q^2)
-\frac{1}{2}\bigl(m_\sigma^2-4M_q^2\bigr)B(m_\sigma^2)
\right],
\label{eq:qm_delta_msigma}
\end{equation}
\begin{equation}
\delta m_\pi^2
=
8 i N_c g^2
\left[
A(M_q^2)
-\frac{1}{2}m_\pi^2 B(m_\pi^2)
\right].
\label{eq:qm_delta_mpi}
\end{equation}

The second on-shell condition fixes the residue of the propagator at
the pole to be unity. For a simple pole this means
\begin{equation}
\mathrm{Res}_{p^2=m_a^2}\, iG_a(p^2) = i,
\end{equation}
which is equivalent to
\begin{equation}
\left.
\frac{dG_a^{-1}(p^2)}{dp^2}
\right|_{p^2=m_a^2}
=1.
\label{eq:qm_OS_residue_condition}
\end{equation}
Substituting the inverse propagator and differentiating gives
\begin{equation}
\left.
\frac{d\Sigma_a^{\mathrm{1PI}}(p^2)}{dp^2}
\right|_{p^2=m_a^2}
+
\delta Z_a
=0,
\end{equation}
and therefore
\begin{equation}
\delta Z_a
=
-
\left.
\frac{d\Sigma_a^{\mathrm{1PI}}(p^2)}{dp^2}
\right|_{p^2=m_a^2}.
\label{eq:qm_delta_Z_general}
\end{equation}

Differentiating the explicit expressions for the self-energies yields
\begin{equation}
\delta Z_\sigma
=
4 i N_c g^2
\left[
B(m_\sigma^2)
+
\bigl(m_\sigma^2-4M_q^2\bigr)B'(m_\sigma^2)
\right],
\label{eq:qm_delta_Zsigma}
\end{equation}
\begin{equation}
\delta Z_\pi
=
4 i N_c g^2
\left[
B(m_\pi^2)
+
m_\pi^2 B'(m_\pi^2)
\right],
\label{eq:qm_delta_Zpi}
\end{equation}
where
\begin{equation}
B'(m_a^2)
\equiv
\left.
\frac{dB(p^2)}{dp^2}
\right|_{p^2=m_a^2}.
\end{equation}

The on-shell conditions therefore determine the mass and
wavefunction counterterms directly in terms of the quark-loop
self-energies. By construction, the renormalized pion and sigma
propagators then have poles at the physical masses \(m_\pi\) and
\(m_\sigma\) and unit residues. These results provide the input needed
to express the counterterms of the original Lagrangian parameters in
the next subsection.

\subsection{Counterterms for the original Lagrangian parameters}

The on-shell conditions determine the counterterms associated with the
physical pion and sigma propagators, namely \(\delta m_\pi^2\),
\(\delta m_\sigma^2\), \(\delta Z_\pi\), and \(\delta Z_\sigma\).
What remains is to relate these quantities to the counterterms of the
parameters that appear in the original Quark--Meson Lagrangian. This
step is necessary because the divergences generated by the quark loops
must ultimately be absorbed into the bare fields and couplings of the
theory \cite{Folkestad2018,Andersen2024}.

After expanding the bare Lagrangian around the vacuum, the shifted
theory may be parametrized in terms of the vacuum expectation value
\(v\), the meson masses \(m_\pi\) and \(m_\sigma\), the Yukawa
coupling \(g\), and the symmetry-breaking parameter \(h\). In this
parametrization the relevant bare quantities are written as
\begin{equation}
m_{\pi,B}^2 = m_\pi^2 + \delta m_\pi^2,
\qquad
m_{\sigma,B}^2 = m_\sigma^2 + \delta m_\sigma^2,
\end{equation}
\begin{equation}
v_B^2 = v^2 + \delta v^2,
\qquad
g_B^2 = g^2 + \delta g^2,
\qquad
h_B = h + \delta h ,
\end{equation}
together with the mesonic wavefunction renormalization constants
\(\delta Z_\pi\) and \(\delta Z_\sigma\) introduced above
\cite{Folkestad2018}. The original parameters \(m^2\) and
\(\lambda\) are then recovered from the tree-level relations
\begin{equation}
m_\sigma^2 = m^2 + \frac{\lambda}{2}v^2,
\qquad
m_\pi^2 = m^2 + \frac{\lambda}{6}v^2,
\qquad
h = m_\pi^2 v ,
\label{eq:qm_tree_relations_shifted}
\end{equation}
which imply
\begin{equation}
m^2 = -\frac{1}{2}\left(m_\sigma^2 - 3m_\pi^2\right),
\qquad
\lambda = 3\,\frac{m_\sigma^2-m_\pi^2}{v^2}.
\label{eq:qm_tree_relations_inverse}
\end{equation}

To one-loop order, products of counterterms may be neglected. The
counterterms for \(m^2\) and \(\lambda\) are therefore obtained by
varying the relations in Eq.~\eqref{eq:qm_tree_relations_inverse}. For
the mass parameter one finds
\begin{equation}
\delta m^2
=
-\frac{1}{2}
\left(
\delta m_\sigma^2
-
3\,\delta m_\pi^2
\right),
\label{eq:qm_delta_m2}
\end{equation}
while the quartic coupling satisfies
\begin{equation}
\delta\lambda
=
3\,\frac{\delta m_\sigma^2-\delta m_\pi^2}{v^2}
-
\lambda
\left(
\frac{\delta v^2}{v^2}
\right).
\label{eq:qm_delta_lambda_prelim}
\end{equation}
The first term comes from the variation of the mass difference
\(m_\sigma^2-m_\pi^2\), while the second term arises because the
coupling is proportional to \(v^{-2}\).

The counterterm \(\delta v^2\) is fixed by the pion wavefunction
renormalization. In the large-\(N_c\) approximation the one-loop
corrections retained here come only from quark loops in mesonic
Green's functions. There is therefore no quark-field renormalization
and no Yukawa-vertex renormalization at this order. The pion field
renormalization is instead absorbed into the vacuum expectation value
and Yukawa coupling so that the constituent quark mass
\(M_q^2=g^2v^2\) is unchanged at this order. This gives
\begin{equation}
\delta g^2 = - g^2 \delta Z_\pi ,
\qquad
\delta v^2 = v^2 \delta Z_\pi .
\label{eq:qm_delta_gv_relation}
\end{equation}
With this identification, Eq.~\eqref{eq:qm_delta_lambda_prelim}
becomes
\begin{equation}
\delta\lambda
=
3\,\frac{\delta m_\sigma^2-\delta m_\pi^2}{v^2}
-
\lambda\,\delta Z_\pi .
\label{eq:qm_delta_lambda}
\end{equation}

The explicit symmetry-breaking parameter is most conveniently fixed
from the linear term in the shifted Lagrangian, for which the tadpole
counterterm contributes together with the variations of
\(m_\pi^2 v\). Using \(\delta t=-8 i N_c g^2 v\,A(M_q^2)\), one finds
\begin{equation}
\delta h
=
-8 i N_c g^2 v\,A(M_q^2)
+
v\,\delta m_\pi^2
+
\frac{1}{2}m_\pi^2 v\,\delta Z_\pi .
\label{eq:qm_delta_h}
\end{equation}

Equations
\eqref{eq:qm_delta_m2},
\eqref{eq:qm_delta_lambda},
\eqref{eq:qm_delta_gv_relation}, and
\eqref{eq:qm_delta_h}
show how the divergences generated by the quark-loop tadpole and
self-energy diagrams are absorbed into redefinitions of the original
parameters of the Quark--Meson Lagrangian. Once the mass and
wavefunction counterterms have been fixed by the on-shell conditions,
the counterterms for \(m^2\), \(\lambda\), \(g\), and \(h\) follow
directly. This completes the renormalization of the vacuum theory at
one-loop order in the large-\(N_c\) approximation.

\subsection{Renormalized vacuum effective potential}

We now return to the vacuum contribution of the fermion determinant
derived in the thermodynamic analysis and show how it is related to
the diagrammatic renormalization carried out above. In the vacuum
limit, the one-loop quark contribution to the effective potential is
the zero-point energy of the fermion field,
\begin{equation}
\Omega_{\mathrm{vac}}^{(1)}
=
-2N_cN_f\,\Lambda^{4-d}
\int\frac{d^{\,d-1}p}{(2\pi)^{d-1}}\,
\sqrt{p^2+M_q^2},
\qquad
M_q = gv ,
\label{eq:qm_Omega_vac_start}
\end{equation}
evaluated in dimensional regularization with \(d=4-2\epsilon\)
\cite{Folkestad2018,Andersen2024}. As discussed earlier, this
integral is ultraviolet divergent.

Using the standard dimensional-regularization formula for the vacuum
integral, one finds
\begin{equation}
\Omega_{\mathrm{vac}}^{(1)}
=
\frac{2N_c\,M_q^4}{(4\pi)^2}
\left[
\frac{1}{\epsilon}
+
\ln\!\left(\frac{\Lambda^2}{M_q^2}\right)
+
\frac{3}{2}
\right]
+\mathcal{O}(\epsilon) .
\label{eq:qm_Omega_vac_div}
\end{equation}
Here \(\Lambda\) denotes the rescaled dimensional-regularization
scale, so the compact logarithmic form is being used.
Since \(M_q=gv\), the divergence is proportional to \(g^4v^4\). It
therefore has the same field dependence as the quartic term in the
mesonic potential and can be absorbed into the counterterms of the
original Lagrangian parameters. This is the renormalizability of the
model in practice: the loop correction does not generate a new
operator, but only renormalizes couplings that are already present
\cite{PeskinSchroeder1995,WeinbergQFT}.

Including the tree-level potential \(U(v)\) and the counterterm
potential \(\delta U(v)\), the vacuum effective potential may be
written as
\begin{equation}
\Omega_{\mathrm{vac}}(v)
=
U(v)
+
\delta U(v)
+
\frac{2N_c\,g^4v^4}{(4\pi)^2}
\left[
\frac{1}{\epsilon}
+
\ln\!\left(\frac{\Lambda^2}{g^2v^2}\right)
+
\frac{3}{2}
\right] .
\label{eq:qm_Omega_vac_with_ct}
\end{equation}
The divergent part of \(\delta U(v)\) is fixed by the counterterms
determined in the previous subsections and cancels the pole term
\(1/\epsilon\) exactly \cite{Folkestad2018}. The renormalized vacuum
potential is therefore finite.

After this cancellation, the remaining vacuum contribution takes the
form
\begin{equation}
\Omega_{\mathrm{vac}}^{\mathrm{ren}}(v)
=
U(v)
+
\delta U_{\mathrm{fin}}(v)
+
\frac{2N_c\,g^4v^4}{(4\pi)^2}
\left[
\ln\!\left(\frac{\Lambda^2}{g^2v^2}\right)
+
\frac{3}{2}
\right],
\label{eq:qm_Omega_vac_ren_generic}
\end{equation}
where \(\delta U_{\mathrm{fin}}(v)\) denotes the finite part of the
counterterm potential. In the on-shell scheme these finite parts are
not arbitrary. They are fixed by the conditions that the vacuum
remains at \(v=f_\pi\) and that the pion and sigma propagators have
their poles at the physical masses with unit residues
\cite{Folkestad2018,Andersen2024}.

This is precisely the vacuum contribution that enters the mean-field
thermodynamic potential. The full zero-temperature potential is
therefore written as
\begin{equation}
\Omega(v;\mu_u,\mu_d)
=
\Omega_{\mathrm{vac}}^{\mathrm{ren}}(v)
+
\Omega_{\mathrm{med}}(v;\mu_u,\mu_d),
\label{eq:qm_full_potential_final}
\end{equation}
where \(\Omega_{\mathrm{med}}\) is the finite medium contribution
derived previously. The renormalization carried out in the vacuum
theory is thus the step that makes the mean-field thermodynamic
potential finite and consistent with the parameter fixing used
throughout the model.

For the applications considered in this thesis, the essential point
is that the ultraviolet divergence of the fermion determinant is fully
absorbed into the counterterms of the vacuum theory. The resulting
renormalized effective potential can then be minimized with respect to
the chiral condensate in exactly the same mean-field approximation as
the one used in the thermodynamic analysis.
